%!TEX root = main.tex
\section{Conclusions}
\label{sec:conclusion}

In this paper, we have presented \castm{}, the first API targeting the spatial-temporal programming model of CGRA architectures for cryptographic workloads. Specifically, for the BabyBear S-box layer on the  OpenEdgeCGRA architecture, \castm{} achieves a \pct{\MetricSboxCastmCycleReductionPct} cycle reduction over hand-optimized CSV (\cc{\MetricSboxCastmCC} vs. \cc{\MetricSboxManualCC}), while reducing source code by approximately \timex{\MetricLocCompressionVsManual}.

RTL validation of the full Poseidon2 hash function on the ultra-low power HEEPsilon SoC gives us some system lessons. The Configuration Register File capacity is the enabling condition: once the kernel fits in the CRF the running times are no longer masked by reload fragmentation. Under that condition,
offloading the full kernel to the CGRA (\textit{CGRA-only}) minimizes the total energy, whereas a hybrid deployment that computes hashes on the CPU and the accelerator (\textit{CPU+CGRA}) offers the highest throughput.

In summary, we have shown that it is possible to enable privacy-preserving verification on wearable biomedical devices through Zero-Knowledge Proofs, on hardware that is both energy-efficient and with reprogramming capability to track the rapid evolution of ZKP systems.

\textbf{Future work.} Two directions follow: (1)~dual-port SRAM or DMA burst modes to reduce the memory stalls that represent now the main source of overhead;
and (2)~evaluation on alternative ZKP fields (Goldilocks, BN254) and post-quantum workloads to test the generality of \castm{}'s API.
